\documentclass[11pt,a4paper]{article}

\usepackage[margin=1in]{geometry}
\usepackage{amsmath,amssymb,amsfonts}
\usepackage{booktabs}
\usepackage{graphicx}
\usepackage{array}
\usepackage{multirow}
\usepackage{hyperref}
\usepackage{microtype}
\usepackage{cite}
\usepackage{authblk}

\input{generated/results_macros.tex}

\title{\textbf{Stability and Generalization of Explainability-Guided Sensor Selection in Subject-Independent sEMG Gesture Recognition}\\[4pt]\large\textit{Preprint intended for Physiological Measurement}}

\author[1]{Minh-Tri Ta}
\author[2,3]{Hoang-Khang Phan}
\author[4,*]{Le-Giang Tran}

\affil[1]{College of Engineering and Computer Science, VinUniversity, Hanoi, Vietnam}
\affil[2]{Department of Biomedical Engineering, Faculty of Applied Science, Ho Chi Minh City University of Technology (HCMUT), Ho Chi Minh City, Vietnam}
\affil[3]{Vietnam National University Ho Chi Minh City, Ho Chi Minh City, Vietnam}
\affil[4]{Faculty of Electronics Technology, Industrial University of Ho Chi Minh City, Ho Chi Minh City, Vietnam}
\affil[*]{Corresponding author: tranlegiang@iuh.edu.vn}

\date{\today}


\begin{document}

\maketitle

\begin{abstract}
Surface electromyography (sEMG) channel reduction is crucial for wearable myoelectric prostheses, yet prior studies rely heavily on user-dependent evaluations, post-hoc threshold heuristics, and unvalidated feature attributions. In this study, we present a rigorous, zero-leakage evaluation framework for explainability-guided sparse sensor selection in subject-independent sEMG gesture recognition. Evaluated on the 43-participant GRABMyo benchmark across 16 active gestures and replicated on NinaPro DB5 (10 subjects), we implement Stability-Regularized Cross-Fitted Attribution (SR-CFA) using TreeSHAP computed strictly across inner training folds. Model performance is assessed under a pre-specified statistical analysis plan ($\delta = 0.02$ Macro-F1 engineering margin) using 10,000 participant-level cluster bootstrap resamples. Full-channel Leave-One-Subject-Out (LOSO) baselines demonstrate that XGBoost (Macro-F1 \FullXGBoostMacroF, 95\% CI [\FullXGBoostMacroFLCB, \FullXGBoostMacroFUCB]) and RBF-SVM (\FullSVMMacroF, 95\% CI [\FullSVMMacroFLCB, \FullSVMMacroFUCB]) are statistically indistinguishable ($p = 0.154$), while significantly outperforming Random Forest (\FullRandomForestMacroF, $p < 0.001$). Under the pre-specified non-inferiority margin, sensor reduction is preserved down to $k^* = 12$ channels ($\text{LCB}_{95\%} = \OptimalKLCB > -0.020$), achieving a 25\% physical sensor reduction without statistical inferiority. Between-subject attribution stability is high across individuals (mean Spearman $\rho = 0.927 \pm 0.038$), concentrating consistently on primary flexor and extensor compartments. A single physically fixed consensus 12-channel layout derived on development subjects generalizes to unseen test subjects with statistical non-inferiority ($\Delta = -0.0124$, 95\% CI $[-0.0179, -0.0077]$). Longitudinal evaluations indicate that multi-day training exposure significantly improves cross-session performance by $+0.0210$ Macro-F1 ($p < 0.001$), demonstrating that multi-session data partially mitigates donning/doffing variability. Replicated on NinaPro DB5, a 12-channel subset yields borderline non-inferiority ($\text{LCB}_{95\%} = -0.0203$). This study provides an audited, leak-free methodological benchmark for population-level wearable neural interface design.
\end{abstract}

\vspace{1em}
\noindent\textbf{Keywords:} Surface electromyography (sEMG), Hand gesture recognition, Explainable AI (XAI), SHAP, Sensor selection, Non-inferiority testing, Leave-one-subject-out (LOSO)

\input{sections/01_introduction.tex}
\input{sections/02_related_work.tex}
\input{sections/03_methods.tex}
\input{sections/04_results.tex}
\input{sections/05_discussion.tex}
\input{sections/06_conclusion.tex}

\section*{Declaration of Competing Interest}
The authors declare no competing financial or non-financial interests.

\section*{Data and Code Availability}
GRABMyo v1.1.0 is available on PhysioNet (\url{https://physionet.org/content/grabmyo/1.1.0/}). NinaPro DB5 is available on the NinaPro repository. Full reproduction scripts, feature tables, and statistical results are open source in the project repository.

\bibliographystyle{IEEEtran}
\bibliography{references}

\end{document}
